\section{Limites et impact élargi}

Les preuves contrôlées se limitent à environ 98,2M paramètres, un tokenizer, un shard FineWeb-Edu, 100M tokens et trois seeds. L'ordre est cohérent, mais l'intervalle à 95\% contient zéro. Il n'établit ni significativité, ni comportement à l'échelle, ni gain downstream, ni supériorité universelle sur GQA.

La comparaison modifie un ensemble : paramétrisation R/W/V et RMSNorm par tête sur R/W. Les six runs n'isolent pas factoriellement la normalisation. Les observations exploratoires à une seed avec écriture lexicale et sans normalisation sont exclues du tableau contrôlé parce qu'elles ne constituent pas des comparaisons appariées sur trois seeds. Des contrôles appariés supplémentaires sont requis avant d'attribuer le gain à un composant.

Le débit dépend de l'implémentation, du logiciel et du H200. Eager est utilisé car la pile compilateur testée produisait des gradients non finis dans un chemin à table lexicale partagée. La pénalité de 2,86\% n'est pas une estimation d'inférence en production. Le cache W/V croît avec le contexte. Prefill long, premier token, distributions de latence, bande passante mémoire et CUDA Graph restent hors portée.

FineWeb-Edu hérite des limites des corpus web : provenance incertaine, déséquilibres et contenus nocifs possibles. L'architecture ne les supprime pas. Le shard est archivé pour reproductibilité exacte et doit être géré selon la licence et la gouvernance de sa source.

Enfin, ce travail est une étude de chemin, non la revendication d'un nouvel opérateur softmax. R et W occupent la frontière technique query/key de SDPA. La terminologie explicite le flux de données ; renommer Q/K/V ne constitue pas seul une contribution algorithmique.